\BeginChapter{06}{Résonance}
\Page{Chapitre 06 / Programme officiel}{Régime sinusoïdal forcé et résonance}{ch06-chapitre-06}
Un excitateur impose sa fréquence ; le système choisit l’amplitude et le déphasage de sa réponse. La résonance est une amplification sélective : sa position dépend de la grandeur observée. Le courant d’un RLC et la tension de sa capacité n’ont donc pas nécessairement leur maximum à la même fréquence.
\begin{loi}{Notions et capacités travaillées}
Impédances de R, L et C ; associations série et parallèle ; représentation complexe ; réponse électrique et mécanique à une excitation sinusoïdale ; acuité et facteur de qualité ; identification à partir de courbes d’amplitude et de phase ; mesure d’une résonance.
\end{loi}

\Sub{Prérequis}
Chapitre 5 : équation canonique, $\omega_0$, Q, extinction du transitoire. Nombres complexes : module, argument et quotient. Dérivation d’un sinus ; amplitude et valeur efficace.

\begin{methode}{La méthode du chapitre}
Fixer la convention temporelle ; construire les impédances ; calculer une amplitude complexe ; prendre module et argument ; étudier les limites ; identifier la grandeur qui résonne ; confronter à une mesure.
\end{methode}

\source{Concordance : annexe 2 de l’arrêté du 5 janvier 2021, rubrique 1.4 / régime sinusoïdal forcé. La puissance moyenne est établie par moyenne temporelle ; aucune théorie de puissance complexe n’est nécessaire.}

\Page{Diagnostic / 15 min}{Avant d’ouvrir le cours}{ch06-diagnostic}
\begin{questions}
\item $s=3\cos(200t+\pi/6)$ V : amplitude, pulsation, fréquence et phase ?
\item Calculer module et argument de $3+4j$.
\item Que vaut $1/j$ ?
\item Dériver $A\cos\omega t$.
\item Quelle fréquence subsiste après extinction du transitoire d’un système linéaire forcé ?
\item Une tension d’amplitude 10 V a quelle valeur efficace ?
\item Rappeler $\omega_0$ et Q du RLC série.
\item Une phase de sortie négative signifie quoi ?
\end{questions}
\begin{docu}{Bilan personnel}
\checkline{Les conversions et les unités sont maîtrisées.}\checkline{Je peux orienter un schéma et annoncer les hypothèses.}\checkline{Je sais expliquer une condition initiale.}Corrigé p.\,\pageref{ch06-corr-diag}.
\end{docu}


\Page{Cours / 1}{Du transitoire au régime établi}{ch06-cours-1}
\Sub{Deux contributions}
Pour un oscillateur stable, la solution complète est la somme d’une solution libre et d’une solution particulière forcée. La première dépend des conditions initiales et s’éteint ; la seconde est à la fréquence de la source. Les calculs de ce chapitre portent sur cette seconde contribution. Ils ne permettent pas de décrire l’allumage à eux seuls.
\begin{loi}{Équation canonique}
\[\ddot s+\frac{\omega_0}{Q}\dot s+\omega_0^2s=\omega_0^2 S_e\cos\omega t.\]Les paramètres $\omega_0$ et Q caractérisent le système ; $\omega$ est choisi par l’excitateur. Le régime établi est $s=S_m\cos(\omega t+\varphi)$.
\end{loi}
\Sub{Attendre avant de mesurer}
En régime sous-amorti, l’enveloppe libre varie comme $e^{-\omega_0t/(2Q)}$. Attendre plusieurs durées $2Q/\omega_0$ après chaque changement de fréquence. Un Q élevé donne un pic étroit mais un établissement lent. Cette observation relie les chapitres 5 et 6.
\begin{piege}{Trois pulsations à distinguer}
La pulsation propre $\omega_0$, la pulsation libre amortie $\omega_d$ et la pulsation du maximum de réponse ne sont pas interchangeables. Nommer systématiquement la grandeur observée.
\end{piege}


\Page{Cours / 2}{La représentation complexe}{ch06-cours-2}
\begin{loi}{Convention unique dans tout le chapitre}
\[s(t)=\Re\{\underline S e^{j\omega t}\},\qquad \underline S=S_m e^{j\varphi}.\]$\underline S$ est une amplitude complexe constante. Dans cette convention, dériver revient à multiplier par $j\omega$ ; intégrer une composante non continue revient à diviser par $j\omega$.
\end{loi}
\Sub{Pourquoi le calcul fonctionne}
Les équations différentielles sont linéaires à coefficients réels. Leur partie réelle est donc une solution réelle. Par exemple $\ddot s$ est représenté par $-\omega^2\underline S$, ce qui transforme une équation différentielle en relation algébrique. Prendre la partie réelle à la fin, après avoir calculé module et phase.
\Sub{Exemple}
Pour $\dot s+s/\tau=e/\tau$, on obtient $(1+j\omega\tau)\underline S=\underline E$. Donc $S_m=E_m/\sqrt{1+(\omega\tau)^2}$ et $\varphi=-\arctan(\omega\tau)$. À $\omega\tau=1$, l’amplitude est divisée par $\sqrt2$ et le retard vaut $45^\circ$.
\begin{methode}{Reconstruire le signal}
Si $\underline S=2-2j$ V, alors $S_m=2\sqrt2$ V et $\varphi=-\pi/4$. Écrire $s(t)=2\sqrt2\cos(\omega t-\pi/4)$ V. Le coefficient réel 2 n’est pas l’amplitude.
\end{methode}


\Page{Cours / 3}{Les impédances élémentaires}{ch06-cours-3}
\begin{loi}{Convention récepteur}
\[\underline U=Z\underline I,\qquad Z_R=R,\quad Z_L=jL\omega,\quad Z_C=\frac1{jC\omega}.\]R, L et C sont idéaux dans ces expressions ; une bobine réelle peut être modélisée par $r+jL\omega$. Chaque Z s’exprime en ohms.
\end{loi}
\Sub{Établissement}
La résistance vérifie $u=Ri$. Pour la bobine $u=L\dot i$, d’où $\underline U=jL\omega\underline I$. Pour la capacité $i=C\dot u$, d’où $\underline I=jC\omega\underline U$. Le courant est en retard de $\pi/2$ sur la tension d’une bobine, et en avance de $\pi/2$ pour un condensateur.
\Sub{Limites physiques}
Quand $\omega$ tend vers zéro, la bobine idéale tend vers un fil et la capacité vers un circuit ouvert. À haute fréquence, les rôles s’inversent dans ce modèle. Ces limites permettent de repérer une erreur de quotient. Elles ne décrivent pas les défauts réels des composants à fréquence arbitrairement élevée.
\begin{docu}{Application}
À $\omega=1000$ rad/s, L = 100 mH donne $Z_L=j100\ \Omega$. C = 10 $\mu\mathrm F$ donne $Z_C=-j100\ \Omega$. En série, ces deux tensions peuvent se compenser malgré des amplitudes non nulles.
\end{docu}


\Page{Cours / 4}{Associer et orienter les dipôles}{ch06-cours-4}
\fig{ch06/rlc-force}{Montage de référence. Le courant traverse tous les dipôles en série ; les tensions de dipôle suivent la convention récepteur.}\begin{loi}{Série et parallèle}
\[Z_s=Z_1+Z_2,\qquad Z_p=\frac{Z_1Z_2}{Z_1+Z_2}.\]Ces formules découlent respectivement de la loi des mailles et de la loi des nœuds. Le diviseur de tension donne $\underline U_2=\underline E Z_2/(Z_1+Z_2)$ uniquement pour les deux impédances réellement en série.
\end{loi}
\Sub{Une charge modifie le circuit}
Un voltmètre ou un étage suivant possède une impédance d’entrée. S’il est en parallèle avec C, remplacer $Z_C$ par $Z_C\parallel Z_e$ avant d’utiliser le diviseur. La résistance interne du GBF appartient aux pertes si $\underline E$ désigne sa f.é.m. à vide ; elle n’est pas à ajouter une seconde fois si l’on utilise la tension mesurée aux bornes du RLC.


\Page{Cours / 5}{Le courant du RLC série}{ch06-cours-5}
\begin{loi}{Réponse en courant}
\[Z=R+j\left(L\omega-\frac1{C\omega}\right),\quad \underline I=\frac{\underline E}{Z}.\]\[I_m=\frac{E_m/R}{\sqrt{1+Q^2(x-1/x)^2}},\quad x=\omega/\omega_0.\]\[\varphi_i=-\arctan[Q(x-1/x)].\]L’argument du dénominateur est dans $]-\pi/2,\pi/2[$ puisque R est positif.
\end{loi}
\fig{ch06/resonance-i}{Maximum du courant à x = 1 pour tout Q positif. La droite pointillée indique le niveau de demi-puissance.}\Sub{Lecture des limites}
À basse fréquence, le condensateur bloque et le courant est en avance ; à haute fréquence, la bobine limite le courant et celui-ci retarde. À $\omega_0$, Z = R, $I_m=E_m/R$ et le courant est en phase avec la source.


\Page{Cours / 6}{La construction de Fresnel}{ch06-cours-6}
\fig{ch06/fresnel}{Addition des tensions complexes dans le cas inductif. Les vecteurs de L et C ont des directions opposées.}\Sub{Une addition vectorielle}
Choisir le courant comme référence de phase ne change aucune tension physique. $\underline U_R$ est réel, $\underline U_L$ vertical vers le haut et $\underline U_C$ vertical vers le bas. Leur somme est la tension source. Le théorème de Pythagore donne directement le module de Z ; l’angle de la somme donne l’avance de la tension sur le courant.
\begin{piege}{Les amplitudes ne s’additionnent pas}
En général $E_m\neq U_{Rm}+U_{Lm}+U_{Cm}$. À la résonance, $U_{Lm}=U_{Cm}=QE_m$ et leurs amplitudes peuvent dépasser $E_m$. Ce sont les amplitudes complexes qui satisfont la loi des mailles.
\end{piege}
\begin{methode}{Contrôle numérique}
R = 20 $\Omega$, L = 0,10 H et C = 10 $\mu\mathrm F$ donnent $\omega_0=1000$ rad/s, Q = 5. Sous $E_m=2$ V, $I_m=0{,}10$ A à la résonance, et les amplitudes sur L et C valent chacune 10 V.
\end{methode}


\Page{Cours / 7}{La résonance en tension}{ch06-cours-7}
\begin{loi}{Tension aux bornes du condensateur}
\[\frac{\underline U_C}{\underline E}=\frac1{1-x^2+jx/Q},\qquad \frac{U_{Cm}}{E_m}=\frac1{\sqrt{(1-x^2)^2+x^2/Q^2}}.\]Le minimum du dénominateur au carré est obtenu pour $x_r^2=1-1/(2Q^2)$. Un maximum à fréquence non nulle existe donc seulement si $Q>1/\sqrt2$.
\end{loi}
\fig{ch06/resonance-c}{Résonance en tension pour Q = 2 ; aucune bosse à fréquence non nulle pour Q = 0,5.}\Sub{Valeur du maximum}
\[\omega_r=\omega_0\sqrt{1-\frac1{2Q^2}},\qquad \frac{U_{Cm,\max}}{E_m}=\frac Q{\sqrt{1-1/(4Q^2)}}.\]À $\omega_0$, le gain vaut exactement Q, sans que ce soit exactement le maximum. Pour Q grand seulement, $\omega_r\simeq\omega_0$ et le gain maximal est proche de Q.


\Page{Cours / 8}{Phase et identification}{ch06-cours-8}
\fig{ch06/phase}{Phase du courant : passage par zéro à la résonance, avance à basse fréquence et retard à haute fréquence.}\Sub{Choisir le bon quadrant}
Pour la capacité, $\varphi_C=-\operatorname{atan2}(x/Q,1-x^2)$ : la phase passe de 0 à $-\pi$, et vaut $-\pi/2$ à $\omega_0$. Employer seulement $-\arctan[(x/Q)/(1-x^2)]$ sans corriger le quadrant donne une erreur de $\pi$ au-dessus de la résonance.
\begin{methode}{Identifier les paramètres}
Le zéro de la phase du courant ou la phase $-\pi/2$ de $u_C$ donnent $\omega_0$. Le gain $U_{Cm}/E_m$ à cette pulsation donne Q. Une bosse en tension donne $\omega_r$, que l’on convertit avec la relation précédente. Vérifier l’accord des deux méthodes.
\end{methode}
\Sub{De la date à la phase}
Si la sortie est retardée d’une durée positive $\Delta t<T/2$, $\varphi_s-\varphi_e=-2\pi\Delta t/T$. La lecture se fait entre deux événements homologues, par exemple deux passages croissants par zéro. La mesure est définie modulo $2\pi$ ; suivre la courbe en fréquence pour dérouler la phase.


\Page{Cours / 9}{Largeur de bande et facteur de qualité}{ch06-cours-9}
\begin{loi}{Bande à demi-puissance du courant}
Les coupures satisfont $I_m=I_{m,\max}/\sqrt2$, donc $|Q(x-1/x)|=1$.\[x_\pm=\sqrt{1+\frac1{4Q^2}}\pm\frac1{2Q},\quad x_-x_+=1.\]\[\Delta\omega=\omega_+-\omega_-=\frac{\omega_0}{Q},\quad Q=\frac{f_0}{f_+-f_-}.\]Ces résultats sont exacts pour le courant du RLC série, avec R constante.
\end{loi}
\Sub{Exemple résolu}
Pour $f_0=1000$ Hz et Q = 5 : $f_-=904{,}99$ Hz, $f_+=1104{,}99$ Hz et $\Delta f=200$ Hz. La moyenne arithmétique des coupures vaut 1004,99 Hz, tandis que leur moyenne géométrique vaut exactement 1000 Hz. Il faut donc éviter de remplacer systématiquement $f_0$ par leur milieu.
\begin{piege}{Le niveau de référence}
Le facteur $1/\sqrt2$ s’applique à une amplitude ; le facteur 1/2 s’applique à sa puissance. Pour la tension sur C, utiliser sans justification la même largeur $\omega_0/Q$ n’est pas exact. À grand Q, les courbes près du pic deviennent voisines.
\end{piege}


\Page{Cours / 10}{L’analogie mécanique}{ch06-cours-10}
\begin{loi}{Masse, ressort et frottement fluide}
\[m\ddot x+b\dot x+kx=F_m\cos\omega t,\quad \omega_0=\sqrt{k/m},\quad Q=\frac{m\omega_0}{b}.\]\[\underline X=\frac{F_m/k}{1-x_r^2+jx_r/Q},\qquad x_r=\omega/\omega_0.\]Le symbole $x_r$ désigne ici une fréquence réduite, pas la position.
\end{loi}
\Sub{Position et vitesse ne résonnent pas pareil}
L’amplitude de position a la même réponse que la tension de C. Elle possède un maximum non nul si $Q>1/\sqrt2$. Comme $\underline V=j\omega\underline X$, la vitesse résonne à $\omega_0$ pour tout Q positif. Son maximum vaut $F_m/b$. La vitesse est alors en phase avec la force tandis que la position retarde de $\pi/2$.
\Sub{Sens physique}
À très basse fréquence, le ressort s’adapte presque statiquement : $X_m\simeq F_m/k$. À très haute fréquence, l’inertie domine : $X_m\simeq F_m/(m\omega^2)$. Le frottement dissipe et limite les amplitudes. Si b = 0 à la résonance exacte, le régime permanent sinusoïdal borné supposé n’existe plus.
\begin{docu}{Correspondances}
$u_C\leftrightarrow x$, $i\leftrightarrow\dot x$, L $\leftrightarrow$ m, R $\leftrightarrow$ b et $1/C\leftrightarrow$ k, à des facteurs d’échelle de source près.
\end{docu}


\Page{Cours / 11}{Énergie reçue et dissipation}{ch06-cours-11}
\Sub{Calcul de la puissance moyenne}
Pour $u=U_m\cos\omega t$ et $i=I_m\cos(\omega t-\theta)$, le produit des cosinus donne\[\langle ui\rangle=\frac{U_m I_m}2\cos\theta.\]Dans le RLC série, cette puissance vaut $RI_m^2/2=RI_{\rm eff}^2$. Les moyennes des puissances reçues par L et C sont nulles en régime établi : leur énergie stockée retrouve sa valeur après une période.
\begin{loi}{Bilan sur une période}
\[\frac{\dd}{\dd t}\left(\frac12Li^2+\frac12Cu_C^2\right)=ei-Ri^2.\]La moyenne du membre de gauche est nulle. À la résonance en courant, $\langle P\rangle=E_m^2/(2R)$, maximum pour une amplitude source fixée.
\end{loi}
\Sub{Interpréter la sélectivité}
À la résonance, les échanges réversibles entre L et C sont importants ; la source compense les pertes. Un Q élevé correspond à de faibles pertes relativement à l’énergie échangée. Dans un montage réel, mesurer R totale et surveiller la surtension : la résistance de la bobine ou du générateur peut limiter Q avant la résistance choisie.
\begin{methode}{Dernier contrôle}
Une puissance moyenne négative reçue par une résistance positive signale une incohérence d’orientation ou de phase. Une puissance moyenne nulle de L n’implique pas que sa puissance instantanée soit toujours nulle.
\end{methode}


\Page{Pratique / TP}{Mesurer une courbe de résonance}{ch06-tp}
\fig{ch06/rlc-mesure}{Placer la résistance de mesure au potentiel de masse pour mesurer $u_R$ sans relier deux points distincts par les masses de l’oscilloscope.}\Sub{Protocole}
Utiliser L = 100 mH, C = 10 $\mu\mathrm F$ et une résistance totale voisine de 20 $\Omega$. Mesurer R totale. Observer simultanément la tension e aux bornes du circuit et $u_R$ sur une résistance accessible de valeur connue $R_m$. Garder $E_m$ constant ; déduire $I_m=U_{Rm}/R_m$. Après stabilisation, relever amplitude et phase pour une quinzaine de fréquences, plus serrées près du pic.
\Sub{Données simulées pour s’entraîner}
Un jeu idéal donne $f_0=159{,}15$ Hz, $f_-=144{,}03$ Hz et $f_+=175{,}86$ Hz. Estimer Q et R totale. Avec des incertitudes-types indépendantes $u(f_0)=0{,}20$ Hz et $u(f_\pm)=0{,}30$ Hz, propager l’incertitude sur Q.
\begin{methode}{Exploitation attendue}
$\Delta f=31{,}83$ Hz, Q = 5,00 et $R=L(2\pi f_0)/Q=20{,}0\ \Omega$. $u(\Delta f)=0{,}424$ Hz et $u(Q)/Q=\sqrt{[u(f_0)/f_0]^2+[u(\Delta f)/\Delta f]^2}$ donnent $u(Q)=0{,}067$. Comparer au Q calculé avec les composants, sans confondre désaccord et erreur certaine.
\end{methode}


\Page{Pratique / Numérique}{Comparer trois résonances}{ch06-numerique}
Le code trace les réponses réduites du courant et de la tension de C pour Q = 0,5 ; 2 ; 5. Les graphes sont calculés sur la même grille ; le tableau donne la position numérique des maxima.
\begin{verbatim}
import numpy as np
import matplotlib.pyplot as plt
x = np.linspace(0.001, 3, 30000)
peaks = []
fig, axes = plt.subplots(2, 1)
for Q in (0.5, 2., 5.):
    gi = 1 / np.sqrt(1 + Q**2*(x-1/x)**2)
    gc = 1 / np.sqrt((1-x*x)**2 + (x/Q)**2)
    peaks.append((Q, x[gi.argmax()], x[gc.argmax()]))
    axes[0].plot(x, gi, label=f"Q={Q}")
    axes[1].plot(x, gc, label=f"Q={Q}")
for ax in axes:
    ax.legend(); ax.grid()
    ax.set_xlabel("omega / omega0")
axes[0].set_ylabel("I / (E/R)")
axes[1].set_ylabel("Uc / E")
plt.tight_layout()
\end{verbatim}
\Sub{Travail et interprétation}
Retrouver $x_i=1$ pour tous les Q. Pour Q = 2 et 5, comparer les maxima en tension à $\sqrt{1-1/(2Q^2)}$, soit 0,9354 et 0,98995. Pour Q = 0,5, le maximum sur la grille est au bord bas : la courbe n’a pas de résonance à fréquence non nulle. Réduire le pas de grille pour distinguer erreur de discrétisation et écart physique.


\Page{Bilan de cours}{Les réflexes à conserver}{ch06-synthese}
\begin{loi}{Les relations et leurs hypothèses}
$Z_R=R$, $Z_L=jL\omega$, $Z_C=1/(jC\omega)$. Pour le RLC série : $\omega_0=1/\sqrt{LC}$, $Q=\sqrt{L/C}/R$. Courant : maximum à $\omega_0$, largeur $\omega_0/Q$. Tension de C : maximum à $\omega_0\sqrt{1-1/(2Q^2)}$ si $Q>1/\sqrt2$.
\end{loi}
\Sub{Une résolution rédigée}
Nommer l’entrée et la sortie ; fixer les orientations ; écrire Z ; contrôler les unités ; calculer module et phase ; vérifier basses et hautes fréquences ; relier le calcul aux courbes. Une identification utilise au moins deux lectures indépendantes.
\begin{piege}{Erreurs à éviter}
Employer une pulsation à la place d’une fréquence ; additionner des modules ; oublier une résistance interne ; imposer une condition initiale à la seule réponse permanente ; confondre résonance de position et de vitesse.
\end{piege}
\begin{docu}{S’autoévaluer}
\checkline{Je retrouve les relations à partir du modèle.}\checkline{Je sais interpréter les limites d’une expression.}\checkline{Je peux dessiner et commenter les courbes.}\checkline{Je relie une mesure à la grandeur calculée.}
\end{docu}


\Page{Exercices / Applications}{Lire un sinus}{ch06-td1}
\source{Énoncés OptiTAB originaux. Les durées sont indicatives ; justifier toute réponse.}
\ex{E01}{Une phase temporelle}{Niveau 1}{10 à 15 min}
$e=4\cos(500t+\pi/3)$ V. Donner $\underline E$, f et la valeur efficace. Écrire l’amplitude complexe de $\dot e$.
\ex{E02}{Revenir au réel}{Niveau 1}{10 à 15 min}
$\underline S=-3+3j$ V. Écrire $s(t)$ à $\omega=100$ rad/s.
\ex{E03}{Mesurer un retard}{Niveau 1}{10 à 15 min}
À f = 200 Hz, une sortie présente un retard de 0,625 ms sur l’entrée. Calculer sa phase relative.


\Page{Exercices / Applications}{Calculer une impédance}{ch06-td2}
\source{Énoncés OptiTAB originaux. Les durées sont indicatives ; justifier toute réponse.}
\ex{E04}{Bobine et condensateur}{Niveau 1}{10 à 15 min}
À f = 50 Hz, calculer Z pour L = 200 mH et pour C = 100 $\mu\mathrm F$.
\ex{E05}{Une série RL}{Niveau 1}{10 à 15 min}
R = 30 $\Omega$, $L\omega=40\ \Omega$ et $E_m=10$ V. Déterminer le courant et sa phase.
\ex{E06}{Un parallèle RC}{Niveau 1}{10 à 15 min}
R = 100 $\Omega$ et $1/(C\omega)=100\ \Omega$. Calculer $Z_R\parallel Z_C$.


\Page{Exercices / Applications}{Déterminer le courant}{ch06-td3}
\source{Énoncés OptiTAB originaux. Les durées sont indicatives ; justifier toute réponse.}
\ex{E07}{La résonance du courant}{Niveau 2}{10 à 15 min}
R = 20 $\Omega$, L = 100 mH, C = 10 $\mu\mathrm F$, $E_m=2$ V. Donner $\omega_0$, f0, Q et le courant maximal.
\ex{E08}{De part et d’autre}{Niveau 2}{10 à 15 min}
Pour les données de E07, calculer $I_m$ et $\varphi_i$ à $\omega=500$ puis 2000 rad/s.
\ex{E09}{Des tensions supérieures à la source}{Niveau 2}{10 à 15 min}
À la résonance de E07, donner les amplitudes des trois tensions. Vérifier la maille complexe.


\Page{Exercices / Méthodes}{Résonance en tension}{ch06-td4}
\source{Énoncés OptiTAB originaux. Les durées sont indicatives ; justifier toute réponse.}
\ex{E10}{Un pic déplacé}{Niveau 2}{10 à 15 min}
Un oscillateur a Q = 2 et $\omega_0=100$ rad/s. Donner la pulsation du maximum en position et le gain maximal réduit.
\ex{E11}{Existe-t-il une résonance ?}{Niveau 2}{10 à 15 min}
Comparer le gain en tension de C pour Q = 0,60 et Q = 0,80. Lequel présente une bosse ?
\ex{E12}{Une tension et sa phase}{Niveau 2}{10 à 15 min}
Pour $x=1$ et Q = 3, écrire $\underline U_C/\underline E$ et interpréter.


\Page{Exercices / Méthodes}{Lire amplitude et phase}{ch06-td5}
\source{Énoncés OptiTAB originaux. Les durées sont indicatives ; justifier toute réponse.}
\ex{E13}{Identifier avec la phase}{Niveau 2}{15 à 25 min}
La phase de $u_C/e$ vaut $-90^\circ$ à 120 Hz, où le gain vaut 4. Déterminer f0, Q et fr.
\ex{E14}{Deux coupures}{Niveau 2}{15 à 25 min}
Le courant atteint son maximum à 1,00 kHz ; ses coupures à demi-puissance sont séparées de 100 Hz. Déterminer Q et leurs valeurs exactes.
\ex{E15}{Passer d’une date à un graphe}{Niveau 2}{15 à 25 min}
À 250 Hz, $u_R$ est en avance de 0,50 ms sur e. Trouver $Q(x-1/x)$.


\Page{Exercices / Méthodes}{Énergie et largeur}{ch06-td6}
\source{Énoncés OptiTAB originaux. Les durées sont indicatives ; justifier toute réponse.}
\ex{E16}{Puissance moyenne}{Niveau 2}{15 à 25 min}
Une source de valeur efficace 5 V alimente une impédance $30+40j\ \Omega$. Calculer la puissance moyenne reçue.
\ex{E17}{La résistance manquante}{Niveau 2}{15 à 25 min}
L = 50 mH, C = 2 $\mu\mathrm F$ et Q mesuré = 5. La résistance ajoutée est 20 $\Omega$. Trouver les autres pertes série.
\ex{E18}{Quand mesurer ?}{Niveau 2}{15 à 25 min}
$\omega_0=1000$ rad/s et Q = 20. Quel temps réduit le transitoire d’un facteur 100 ?


\Page{Exercices / Approfondissement}{Passer à la mécanique}{ch06-td7}
\source{Énoncés OptiTAB originaux. Les durées sont indicatives ; justifier toute réponse.}
\ex{E19}{Une masse forcée}{Niveau 3}{15 à 25 min}
m = 0,20 kg, k = 20 N/m, b = 0,40 kg/s, $F_m=0{,}10$ N. Calculer Q, $\omega_0$ et les amplitudes à $\omega_0$.
\ex{E20}{Dissipation mécanique}{Niveau 3}{15 à 25 min}
Pour E19 à $\omega_0$, calculer la puissance moyenne et vérifier le bilan.
\ex{E21}{Borne sur Q}{Niveau 3}{15 à 25 min}
On impose $\omega_r\geq0{,}99\omega_0$ pour la position. Déterminer la borne minimale de Q.


\Page{Exercices / Approfondissement}{Discuter un modèle}{ch06-td8}
\source{Énoncés OptiTAB originaux. Les durées sont indicatives ; justifier toute réponse.}
\ex{E22}{Une géométrie de coupures}{Niveau 3}{15 à 25 min}
Démontrer $\omega_-\omega_+=\omega_0^2$, puis retrouver $\omega_0$ pour des coupures à 900 et 1100 rad/s.
\ex{E23}{Un capteur charge la capacité}{Niveau 3}{15 à 25 min}
Une résistance d’entrée $R_e$ est placée en parallèle sur C. Écrire l’impédance de sortie, puis le transfert pour un RLC série.
\ex{E24}{Une résonance sans pertes ?}{Niveau 3}{15 à 25 min}
Pour b = 0 et $\omega=\omega_0$, pourquoi le transfert forcé diverge-t-il ? Quelle conclusion physique faut-il tirer ?


\Page{Défi avancé / D1}{Concevoir une résonance sélective}{ch06-defi1}
\label{ch06-ex-D1}\source{Problème OptiTAB original. Niveau 3 ; 40 à 60 minutes. Indices p.\,\pageref{ch06-indices} ; corrigé p.\,\pageref{ch06-corr-defi1}.}
Un RLC série doit résonner en courant à 1,00 kHz avec une largeur de bande de 100 Hz. L = 50,0 mH ; la bobine possède une résistance de 8,0 $\Omega$. La source appliquée au RLC a une amplitude de 1,00 V.\begin{questions}
\item Déterminer C, Q et R totale. En déduire la résistance externe.
\item Calculer les deux fréquences de coupure et la phase du courant en chacune.
\item Calculer le courant maximal, la surtension sur C à f0 et la puissance moyenne.
\item Un capteur ne supporte que 8 V d’amplitude sur C. Peut-on conserver simultanément les caractéristiques imposées ? Proposer une modification de la source.
\item Estimer le délai d’attente pour réduire le transitoire à 1 \%.
\end{questions}


\Page{Défi avancé / D2}{Identifier sans connaître les composants}{ch06-defi2}
\label{ch06-ex-D2}\source{Problème OptiTAB original. Niveau 3 ; 40 à 60 minutes. Indices p.\,\pageref{ch06-indices} ; corrigé p.\,\pageref{ch06-corr-defi2}.}
Un oscillateur mécanique de masse connue m = 0,20 kg reçoit une force sinusoïdale d’amplitude $F_m=0{,}10$ N. La vitesse maximale vaut 0,25 m/s à $\omega=10$ rad/s. Les coupures à demi-puissance sont distantes de 2 rad/s.\begin{questions}
\item Déterminer b, k et Q de deux manières cohérentes.
\item Prédire la position du maximum de déplacement et son amplitude.
\item Calculer la puissance moyenne reçue à la résonance de vitesse.
\item À $\omega_0$, un relevé annonce que la position est en avance de $90^\circ$ sur la force. Le relevé est-il cohérent ? Quels branchements ou définitions contrôler ?
\end{questions}


\Page{Défi avancé / D3}{Les pertes du générateur changent la lecture}{ch06-defi3}
\label{ch06-ex-D3}\source{Problème OptiTAB original. Niveau 3 ; 40 à 60 minutes. Indices p.\,\pageref{ch06-indices} ; corrigé p.\,\pageref{ch06-corr-defi3}.}
Un générateur de f.é.m. $E_g$ et de résistance r = 50 $\Omega$ alimente R = 20 $\Omega$, L = 100 mH et C = 10 $\mu\mathrm F$ en série. On peut soit rapporter $u_C$ à $E_g$, soit à la tension e réellement appliquée au groupe RLC.\begin{questions}
\item Établir les deux transferts et calculer leur Q respectif.
\item À $\omega_0$, comparer les gains et la tension réellement appliquée au RLC.
\item Pour chaque transfert, existe-t-il une résonance en tension à fréquence non nulle ?
\item Expliquer pourquoi garder la consigne du GBF constante ne suffit pas à garantir e constante pendant le balayage.
\end{questions}


\Page{Évaluation / DM}{DM - Un résonateur à calibrer}{ch06-dm}
\source{Devoir maison OptiTAB original ; 1 h 30 à 2 h ; calculatrice autorisée. Barème pédagogique sur 20.}
Un RLC série possède L = 80,0 mH, C inconnue et R totale inconnue. Sous une tension d’amplitude constante 2,00 V, on mesure $f_0=500$ Hz et un courant maximal de 40,0 mA. Les lectures à demi-puissance donnent $\Delta f=99{,}5$ Hz.\Sub{A. Identification - 8 points}
\begin{questions}
\item Établir le transfert en courant et sa phase. (3)
\item Déterminer C, R totale et Q avec les composants. (3)
\item Comparer au Q donné par la largeur. (2)
\end{questions}

\Sub{B. Tension et énergie - 8 points}
\begin{questions}
\item Calculer l’amplitude sur C à f0 et celle du maximum véritable. (3)
\item Donner la fréquence de ce maximum. (2)
\item Établir le bilan énergétique moyen et calculer la puissance reçue à f0. (3)
\end{questions}

\Sub{C. Mesure - 4 points}
\begin{questions}
\item Préciser les grandeurs à maintenir constantes, les précautions de masse et le temps de stabilisation à 1 \%. (4)
\end{questions}



\Page{Évaluation / DS / 1 sur 2}{DS - Deux oscillateurs, une méthode}{ch06-ds1}
\source{Devoir surveillé OptiTAB original ; durée 2 h ; calculatrice autorisée ; note sur 20. Les deux pages forment un seul sujet.}
\Sub{A. Représentation complexe - 6 points}
On impose $e=3\cos\omega t$ V à un RLC série : R = 30 $\Omega$, L = 60 mH, C = 6,0 $\mu\mathrm F$.\begin{questions}
\item Établir les trois impédances et le courant complexe. (3)
\item Donner $\omega_0$, Q, f0 et le courant maximal. (3)
\end{questions}

\Sub{B. Résonance et phase - 8 points}
\begin{questions}
\item Calculer les coupures exactes du courant et leur phase. (3)
\item Calculer la tension de C et sa phase à f0. (2)
\item Le maximum de cette tension se produit-il à f0 ? Calculer sa fréquence et son amplitude. (3)
\end{questions}



\Page{Évaluation / DS / 2 sur 2}{Approfondir et conclure}{ch06-ds2}
\Sub{C. Une expérience mécanique - 6 points}
Un système masse-ressort présente m = 0,30 kg, k = 30 N/m et b = 0,60 kg/s. On applique $F=0{,}12\cos\omega t$ N.\begin{questions}
\item Déterminer $\omega_0$, Q et le transfert en déplacement. (2)
\item À $\omega_0$, déterminer les amplitudes de position et vitesse et leurs phases relatives à F. (2)
\item Calculer la puissance moyenne et expliquer le rôle du frottement dans le régime établi. (2)
\end{questions}

\begin{docu}{Gestion du temps}
Lecture : 10 min ; parties A et B : 70 min ; partie C : 30 min ; vérification : 10 min. Une expression littérale, une unité et une interprétation accompagnent les résultats.
\end{docu}


\Page{Aides progressives}{Des indices pour repartir}{ch06-indices}
\Sub{Défi D1}
1. La largeur fixe Q. 2. R totale = L\ensuremath{\omega_0}/Q. 3. La tension sur C à f0 vaut QE. 4. Modifier l’amplitude ne change pas les paramètres d’un modèle linéaire.
\Sub{Défi D2}
1. Vmax = Fm/b. 2. \ensuremath{\Delta}\ensuremath{\omega} = b/m. 3. Le déplacement suit le transfert passe-bas du second ordre.
\Sub{Défi D3}
1. Écrire la maille en incluant r seulement quand l’entrée est Eg. 2. e = ZRLC I. 3. Une tension mesurée permet de normaliser le transfert.
\Sub{Exercices}
E06 : utiliser l’admittance. E10 : minimiser le dénominateur au carré. E14 : utiliser les coupures exactes, pas une symétrie arithmétique. E23 : remplacer d’abord la branche C par son parallèle.
\Sub{Évaluations}
DM : la résistance vient directement de la résonance en courant. DS : conserver $\omega_0$ en rad/s jusqu’à la conversion finale ; une dérivation multiplie l’amplitude par $\omega$ et ajoute $\pi/2$ à la phase.


\Page{Corrigés / Diagnostic}{Réactiver les bons réflexes}{ch06-corr-diag}
\Sub{D1}
$S_m=3$ V ; $\omega=200$ rad/s ; $f=100/\pi=31{,}83$ Hz ; $\varphi=\pi/6$.

\Sub{D2}
Module 5 ; argument $\arctan(4/3)=53{,}13^\circ$.

\Sub{D3}
$1/j=-j$ car $j^2=-1$.

\Sub{D4}
$-A\omega\sin\omega t=A\omega\cos(\omega t+\pi/2)$.

\Sub{D5}
Celle de la source, si l’entrée ne contient qu’une fréquence et si le régime stable existe.

\Sub{D6}
$10/\sqrt2=7{,}071$ V pour une sinusoïde sans composante continue.

\Sub{D7}
$\omega_0=1/\sqrt{LC}$ et $Q=L\omega_0/R=\sqrt{L/C}/R$.

\Sub{D8}
La sortie retarde sur l’entrée, avec la convention $\cos(\omega t+\varphi)$ et une comparaison à fréquence commune.


\Page{Corrigés / Exercices}{Lire un sinus}{ch06-corr-td1}
\sol{E01}{Une phase temporelle}
$\underline E=4e^{j\pi/3}=2+2\sqrt3j$ V ; $f=500/(2\pi)=79{,}58$ Hz ; $E_{\rm eff}=2\sqrt2$ V. La dérivée a pour amplitude $j500\underline E$, de module 2000 V/s et de phase $5\pi/6$.
\sol{E02}{Revenir au réel}
Module $3\sqrt2$ V, argument $3\pi/4$ dans le deuxième quadrant. Donc $s=3\sqrt2\cos(100t+3\pi/4)$ V. La tangente seule ne distingue pas les quadrants.
\sol{E03}{Mesurer un retard}
T = 5 ms ; $\varphi=-2\pi(0{,}625/5)=-\pi/4$. Le signe négatif traduit le retard dans la convention du cours.


\Page{Corrigés / Exercices}{Calculer une impédance}{ch06-corr-td2}
\sol{E04}{Bobine et condensateur}
$Z_L=j20\pi=j62{,}83\ \Omega$ ; $Z_C=-j100/\pi=-j31{,}83\ \Omega$. Les deux modules ne sont pas égaux.
\sol{E05}{Une série RL}
$Z=30+40j\ \Omega$, $|Z|=50\ \Omega$. $I_m=0{,}20$ A, $\varphi_i=-\arctan(4/3)=-53{,}13^\circ$. La valeur efficace est 0,1414 A.
\sol{E06}{Un parallèle RC}
$Y=1/R+jC\omega=0{,}01(1+j)$ S, donc $Z=50-50j\ \Omega$, de module $50\sqrt2\ \Omega$. Une association parallèle se traite commodément avec les admittances.


\Page{Corrigés / Exercices}{Déterminer le courant}{ch06-corr-td3}
\sol{E07}{La résonance du courant}
$\omega_0=1000$ rad/s, $f_0=159{,}15$ Hz, Q = 5 et $I_{m,\max}=0{,}10$ A. Le maximum est obtenu lorsque la partie imaginaire de Z s’annule.
\sol{E08}{De part et d’autre}
Les fréquences réduites sont 0,5 et 2, avec $Q(x-1/x)=\mp7{,}5$. Dans les deux cas $I_m=0{,}1/\sqrt{57{,}25}=13{,}22$ mA. Les phases valent respectivement $+82{,}41^\circ$ et $-82{,}41^\circ$.
\sol{E09}{Des tensions supérieures à la source}
$U_{Rm}=2$ V ; $U_{Lm}=U_{Cm}=10$ V. Avec le courant réel : $\underline U_R=2$, $\underline U_L=10j$, $\underline U_C=-10j$ V ; leur somme est bien 2 V.


\Page{Corrigés / Exercices}{Résonance en tension}{ch06-corr-td4}
\sol{E10}{Un pic déplacé}
$\omega_r=100\sqrt{7/8}=93{,}54$ rad/s. Le gain maximal est $2/\sqrt{15/16}=8/\sqrt{15}=2{,}066$. À $\omega_0$, le gain n’est que 2.
\sol{E11}{Existe-t-il une résonance ?}
Le seuil est $1/\sqrt2\simeq0{,}7071$. Q = 0,60 : décroissance depuis la valeur statique. Q = 0,80 : maximum non nul à $x_r=\sqrt{7/32}=0{,}4677$. Un régime libre oscillant, possible dès Q > 0,5, ne garantit donc pas cette bosse.
\sol{E12}{Une tension et sa phase}
$H=1/(j/3)=-3j$. Gain 3 et phase $-\pi/2$. La tension de C retarde d’un quart de période ; le courant est en phase avec la source.


\Page{Corrigés / Exercices}{Lire amplitude et phase}{ch06-corr-td5}
\sol{E13}{Identifier avec la phase}
$f_0=120$ Hz, Q = 4. $f_r=120\sqrt{1-1/32}=118{,}11$ Hz. Ces conclusions supposent la forme canonique du RLC et l’absence de charge parasite importante.
\sol{E14}{Deux coupures}
Q = 10. $f_\pm=1000[\sqrt{1+1/400}\pm1/20]$, soit 951,25 et 1051,25 Hz. Leur produit vaut $f_0^2$.
\sol{E15}{Passer d’une date à un graphe}
$\varphi_i=+\pi/4$. Puisque $\varphi_i=-\arctan[Q(x-1/x)]$, on a $Q(x-1/x)=-1$. C’est la coupure basse du courant.


\Page{Corrigés / Exercices}{Énergie et largeur}{ch06-corr-td6}
\sol{E16}{Puissance moyenne}
$I_{\rm eff}=5/50=0{,}10$ A et $\cos\theta=30/50=0{,}60$. $P=5\times0{,}10\times0{,}60=0{,}30$ W, également $RI_{\rm eff}^2$.
\sol{E17}{La résistance manquante}
$\sqrt{L/C}=158{,}11\ \Omega$ ; R totale = 31,62 $\Omega$. Les pertes non comptées valent donc 11,62 $\Omega$. Si la tension d’entrée est mesurée sur le seul RLC, la résistance interne du GBF est exclue de ce R totale.
\sol{E18}{Quand mesurer ?}
La constante d’enveloppe vaut $2Q/\omega_0=0{,}040$ s. Il faut $t=0{,}040\ln100=0{,}1842$ s. La consigne concerne l’enveloppe du mode libre sous-amorti.


\Page{Corrigés / Exercices}{Passer à la mécanique}{ch06-corr-td7}
\sol{E19}{Une masse forcée}
$\omega_0=10$ rad/s et Q = 5. $X_m=F_m/(b\omega_0)=0{,}025$ m ; $V_m=F_m/b=0{,}25$ m/s. La position retarde de $\pi/2$, la vitesse est en phase avec la force.
\sol{E20}{Dissipation mécanique}
$\langle P_f\rangle=bV_m^2/2=0{,}0125$ W. La force et la vitesse étant en phase, $\langle Fv\rangle=F_mV_m/2=0{,}0125$ W. L’énergie mécanique moyenne ne varie plus.
\sol{E21}{Borne sur Q}
$1-1/(2Q^2)\geq0{,}99^2$ donne $Q\geq[2(1-0{,}9801)]^{-1/2}=5{,}013$. La présence du pic et l’approximation de sa position sont deux exigences distinctes.


\Page{Corrigés / Exercices}{Discuter un modèle}{ch06-corr-td8}
\sol{E22}{Une géométrie de coupures}
Le produit des deux racines réduites $\sqrt{1+1/(4Q^2)}\pm1/(2Q)$ est 1. Donc $\omega_0=\sqrt{900\times1100}=994{,}99$ rad/s et Q = 994,99/200 = 4,975.
\sol{E23}{Un capteur charge la capacité}
$Z_p=[1/R_e+jC\omega]^{-1}$. $H=Z_p/[R+jL\omega+Z_p]$. En développant : $H=1/[1+R/R_e-LC\omega^2+j\omega(RC+L/R_e)]$. L’ancien Q ne s’utilise plus sans recalcul.
\sol{E24}{Une résonance sans pertes ?}
L’équation $\ddot x+\omega_0^2x=(F_m/m)\cos\omega_0t$ possède le terme particulier $x_p=F_mt\sin(\omega_0t)/(2m\omega_0)$. Son amplitude croît avec t. Il n’existe pas de régime permanent borné à cette fréquence dans ce modèle idéal ; la divergence n’est pas une amplitude réelle infinie établie instantanément.


\Page{Corrigé du défi D1}{Concevoir une résonance sélective}{ch06-corr-defi1}
C = $1/[L(2\pi f_0)^2]=0{,}5066\ \mu\mathrm F$, Q = 10 et R totale = 31,416 $\Omega$. Ajouter 23,416 $\Omega$. Les coupures sont 951,25 et 1051,25 Hz ; les phases du courant sont $+45^\circ$ et $-45^\circ$. Le courant maximal vaut 31,83 mA et $U_{Cm}(f_0)=10$ V. La puissance moyenne est $1/(2R)=15{,}92$ mW. Le maximum en tension est même légèrement supérieur à 10 V : gain $10/\sqrt{1-1/400}=10{,}0125$. Une limitation à 8 V sur toute la courbe impose $E_m\leq0{,}7990$ V. L et C et les résistances peuvent rester inchangés. La durée d’enveloppe vaut $2Q/\omega_0=3{,}183$ ms ; attendre au moins 14,66 ms pour un facteur 100, puis tenir compte des délais instrumentaux. Un cahier des charges doit préciser la grandeur limitée : courant, tension ou puissance.

\Page{Corrigé du défi D2}{Identifier sans connaître les composants}{ch06-corr-defi2}
$b=F_m/V_{\max}=0{,}40$ kg/s. $k=m\omega_0^2=20$ N/m. $Q=m\omega_0/b=5$ ; la largeur donne aussi Q = 10/2 = 5 et b = m\ensuremath{\Delta}\ensuremath{\omega} = 0,40 kg/s. Le maximum de déplacement est à $\omega_r=10\sqrt{0{,}98}=9{,}8995$ rad/s. Son amplitude vaut $(F_m/k)Q/\sqrt{1-1/(4Q^2)}=25{,}126$ mm. La puissance moyenne vaut $bV_{\max}^2/2=12{,}5$ mW. La position doit retarder de $90^\circ$ sur la force. Un capteur de déplacement inversé ajoute $\pi$ à la phase et ferait apparaître une avance de $90^\circ$. Vérifier aussi la polarité de la mesure de force et l’ordre des voies. L’accord des modules ne suffit pas à valider la phase.

\Page{Corrigé du défi D3}{Les pertes du générateur changent la lecture}{ch06-corr-defi3}
Par rapport à $E_g$, $H_g=1/[1-LC\omega^2+j(R+r)C\omega]$ et $Q_g=100/70=1{,}4286$. Par rapport à e, $H_e=1/[1-LC\omega^2+jRC\omega]$ et $Q_e=5$. À $\omega_0=1000$ rad/s, $|u_C/E_g|=1{,}4286$, $|u_C/e|=5$ et $e/E_g=R/(R+r)=2/7$. Les deux Q dépassent $1/\sqrt2$ : les deux transferts ont une bosse, respectivement à $x_r=\sqrt{1-1/(2Q_g^2)}=0{,}8689$ et 0,98995. La charge du générateur dépend de la fréquence, donc sa chute interne varie. Le rapport $u_C/e$ reste celui du RLC même si e varie ; une courbe de $u_C$ seule avec consigne fixe décrit en revanche tout le montage. Il faut annoncer précisément le port d’entrée utilisé.

\Page{Corrigé / DM / Sur 20}{DM - Un résonateur à calibrer}{ch06-corr-dm}
\Sub{A - 8 points}
$\underline I=\underline E/[R+j(L\omega-1/(C\omega))]$, phase opposée à l’argument du dénominateur (3). R = E/Imax = 50,0 $\Omega$, C = $1/[0{,}080(1000\pi)^2]=1{,}2665\ \mu\mathrm F$, Q = L\ensuremath{\omega_0}/R = 5,0265 (3). La largeur donne 500/99,5 = 5,0251, accord avec l’arrondi des données (2).
\Sub{B - 8 points}
À f0, $U_{Cm}=QE_m=10{,}053$ V ; maximum $QE_m/\sqrt{1-1/(4Q^2)}=10{,}103$ V (3). $f_r=f_0\sqrt{1-1/(2Q^2)}=495{,}03$ Hz (2). La variation moyenne des énergies stockées est nulle ; $\langle ei\rangle=RI_m^2/2=0{,}0400$ W (3).
\Sub{C - 4 points}
Mesurer e aux bornes du RLC et maintenir son amplitude, ou normaliser chaque point (1). Placer les masses au même nœud et mesurer le courant par une résistance accessible (1). La constante d’enveloppe vaut $2L/R=3{,}20$ ms ; un facteur 100 demande $3{,}20\ln100=14{,}74$ ms (2). Le barème total est 20.


\Page{Corrigé / DS / 1 sur 2}{Calculer et interpréter}{ch06-corr-ds1}
\Sub{A - 6 points}
$Z_R=R$, $Z_L=jL\omega$, $Z_C=1/(jC\omega)$ ; $\underline I=3/[30+j(0{,}060\omega-1/(6\times10^{-6}\omega))]$ (3). $\omega_0=1666{,}67$ rad/s, Q = 10/3, $f_0=265{,}26$ Hz et $I_{m,\max}=0{,}100$ A (3).
\Sub{B - 8 points}
$x_\pm=\sqrt{1+0{,}15^2}\pm0{,}15$, donc $\omega_-=1435{,}31$ et $\omega_+=1935{,}31$ rad/s ; phases $+\pi/4$ et $-\pi/4$ ; largeur 500 rad/s (3). À f0, $U_{Cm}=10{,}0$ V et phase $-\pi/2$ (2). Q dépasse $1/\sqrt2$ ; $\omega_r=1666{,}67\sqrt{0{,}955}=1628{,}73$ rad/s, $f_r=259{,}22$ Hz. Le gain maximal vaut 3,3715, donc $U_{Cm,\max}=10{,}114$ V (3).


\Page{Corrigé / DS / 2 sur 2}{Calculer et interpréter}{ch06-corr-ds2}
\Sub{C - 6 points}
$\omega_0=10$ rad/s ; Q = 5 ; $\underline X=(0{,}12/30)/[1-(\omega/10)^2+j\omega/50]$ m (2). À $\omega_0$, $X_m=0{,}020$ m avec retard $\pi/2$, $V_m=0{,}20$ m/s en phase avec F (2). $P=bV_m^2/2=0{,}012$ W. La source compense exactement la dissipation moyenne, sans augmentation moyenne d’énergie stockée (2).
\begin{methode}{Lire sa copie}
Une erreur sur Q se répercute sur la largeur, la surtension et l’établissement. Refaire d’abord l’identification des coefficients de l’équation. Ne pas retirer deux fois les points d’une même erreur de calcul dans les questions dépendantes.
\end{methode}


\Page{Support authentique / 1}{TP publié - Résonance du RLC}{ch06-sources-1}
\source{Source consultée le 25 octobre 2026. MPSI, TP12, année 2025-2026, p. 1-2. Sélection : modèle et étude rapide.}
\begin{loi}{Accéder au document d’origine}
\href{https://mpsi.schleck.ovh/fichiers/tp/tp12_rlc_rsf.pdf}{Ouvrir le sujet original et conserver ses figures et données.}
\end{loi}
\Sub{Parcours de travail}
Refaire la préparation ; mesurer le courant à partir de la résistance. Examiner le rôle du suiveur et de la résistance de bobine. Comparer deux résistances.
\Sub{Repères de correction indépendants}
Le gain maximal de la tension résistive est $R/(R+r)$, tandis que Q dépend de R + r. À résonance, la phase résistive est nulle. La largeur du courant vérifie $\Delta\omega=(R+r)/L$ lorsque l’entrée est mesurée sur le groupe série considéré.
\source{Ce parcours ne reproduit pas l’épreuve complète. Les questions non ciblées peuvent demander des chapitres ultérieurs. Le temps de travail proposé concerne uniquement cette sélection.}

\Page{Support authentique / 2}{TD publié - Une résonance sélective}{ch06-sources-2}
\source{Source consultée le 25 octobre 2026. MPSI, TD7, 2025-2026, exercice 3, p. 1.}
\begin{loi}{Accéder au document d’origine}
\href{https://mpsi.schleck.ovh/fichiers/td/td7_filtrage_lineaire.pdf}{Ouvrir le sujet original et conserver ses figures et données.}
\end{loi}
\Sub{Parcours de travail}
Traiter les questions 1 à 6 à partir de la fonction fournie ; distinguer gain, gain en décibels et puissance.
\Sub{Repères de correction indépendants}
Le gain est $[1+Q^2(x-1/x)^2]^{-1/2}$. Le maximum vaut 1 à x = 1. Les coupures satisfont $Q|x-1/x|=1$ ; leur différence est $1/Q$ en fréquence réduite. Les pentes asymptotiques sont +20 puis -20 dB par décade.
\source{Ce parcours ne reproduit pas l’épreuve complète. Les questions non ciblées peuvent demander des chapitres ultérieurs. Le temps de travail proposé concerne uniquement cette sélection.}

\Page{Références et suivi}{Préparer la suite}{ch06-references}
\Sub{Référence réglementaire}
\href{https://www.education.gouv.fr/bo/21/Special1/ESRS2035779A.htm}{Bulletin officiel spécial n° 1 du 11 février 2021, arrêté du 5 janvier 2021.}\par
\href{https://cache.media.education.gouv.fr/file/SPE1-MEN-MESRI-4-2-2021/64/8/spe779_annexe_1373648.pdf}{Annexe officielle regroupée}, rubrique 1.4 / régime sinusoïdal forcé. La puissance moyenne est établie par moyenne temporelle ; aucune théorie de puissance complexe n’est nécessaire.
\Sub{Nature des contenus}
Le cours, les 24 exercices, les trois défis et les évaluations OptiTAB sont rédigés pour cette collection. Les graphiques sont calculés avec les équations annoncées. Les supports authentiques sont identifiés sur les pages précédentes ; leurs énoncés et figures restent dans le document d’origine.
\Sub{Compétences à confirmer}
\checkline{Je peux retrouver l’équation sans mémoriser une suite de signes.}
\checkline{Je vérifie les conditions initiales et les limites.}
\checkline{Je peux confronter un calcul, un graphe et un bilan énergétique.}
\checkline{Je sais corriger mon DS et refaire une question sans aide.}
\begin{methode}{Prochaine étape}
Le chapitre 7 applique ces réponses fréquentielles au traitement de signaux contenant plusieurs harmoniques.
\end{methode}

